% =============================================================
\section{Reproducibility}\label{app:reproducibility}
% =============================================================

The build pipeline binds the entire compute environment, not just the
Python packages. A \texttt{requirements.txt} inside a Dockerfile pins
Python packages and a base OS image but leaves the C and
C\nolinkurl{++} compilers, the CUDA toolchain, the Python interpreter
build, and the cloud topology unspecified --- each of which can drift
between machines and change numerical results. The stack we use
closes those gaps.

\paragraph{Python dependencies (\texttt{uv}).} A Rust-implemented
package manager that resolves direct and transitive Python
dependencies and records every resulting wheel and its hash in a
lockfile (\texttt{uv.lock}). Two machines sharing the lockfile
install the same versions of every Python package, regardless of
install order.

\paragraph{System environment (Nix).} A Nix flake (\texttt{flake.nix})
specifies the Python interpreter (3.11), CUDA runtime, JAX, the C
compiler that built them, and command-line tooling. Store paths are
determined by the hash of the full transitive build graph (pinned via
\texttt{flake.lock}), so the same flake on the same target system
produces the same dependency closure.

\paragraph{Cloud infrastructure (OpenTofu).} The TPU pod-slice, Cloud
SQL Postgres database, and GCS artifact bucket are declared as
OpenTofu configuration in \texttt{infra/tofu/}; \texttt{tofu apply}
brings the topology up, \texttt{tofu destroy} tears it down.

\paragraph{Reproducing the figures.} A reader entering the Nix shell
gets a deterministic Python and CUDA environment; the project
\texttt{README} carries the exact commands. Reproducing the figures
from pinned data does not require any TPU-scale workload.

\paragraph{Artifact set.} The curated data backing every figure ---
parity \texttt{.npz} files, the strength--duration and biphasic-grid
Zarr stores, the throughput sweeps, and the gradient receipt ---
totals a few megabytes and is published as a separate artifact
distribution rather than committed to the source repository, available
at \url{https://github.com/Marco-Christiani/jaxley-extracellular/releases/tag/paper-v0.1}.
